\chapter*{Introduzione}
\rhead[\fancyplain{}{\bfseries
INTRODUZIONE}]{\fancyplain{}{\bfseries\thepage}}
\lhead[\fancyplain{}{\bfseries\thepage}]{\fancyplain{}{\bfseries
INTRODUZIONE}}
\addcontentsline{toc}{chapter}{Introduzione}

Obbiettivo della tesi: analizzare lo stato dell'arte

\section{Contesto e obiettivi della tesi}

Contesto: perch\'e il PM \`e centrale nei progetti IT
Obiettivo della tesi esplicitato: analizzare lo stato dell'arte del project management, con un focus sulle metodologie tradizionali e Agile,
e confrontare le due approcci in termini di efficacia e adozione nelle aziende.

\section{Metodologia di analisi}

La tesi \`e una revisione della letteratura basata su fonti ufficiali e ricerche accademiche.
Per le fonti (standard ufficiali, paper accademici, report di settore) sono stati definiti criteri di selezione basati su: rilevanza, autorevolezza e aggiornamento.

\section{Struttura del documento}

La struttura del documento \`e la seguente: metodi tradizionali -> Agile -> confronto -> stato delle aziende.
